\section{Phân tích yêu cầu phi chức năng}

Yêu cầu phi chức năng mô tả các tiêu chí chất lượng mà hệ thống cần đáp ứng trong quá trình vận hành. Khác với yêu cầu chức năng tập trung vào việc hệ thống “làm gì”, yêu cầu phi chức năng tập trung vào việc hệ thống cần hoạt động “như thế nào”, bao gồm độ chính xác, khả năng truy vết, khả năng mở rộng, khả năng vận hành, bảo mật và phân quyền.

Đối với hệ thống KMA Chatbot tuyển sinh, yêu cầu phi chức năng có vai trò đặc biệt quan trọng vì hệ thống xử lý các thông tin liên quan đến tuyển sinh, học phí, ngành đào tạo, thủ tục, văn bản và tài liệu chính thức. Nếu câu trả lời không đúng nguồn hoặc dữ liệu không được kiểm soát, người dùng có thể hiểu sai thông tin tuyển sinh. Vì vậy, hệ thống cần được thiết kế để đảm bảo câu trả lời có căn cứ, dữ liệu có thể kiểm tra lại và các chức năng quản trị được phân quyền rõ ràng.

\begin{figure}[H]
    \centering
    \includegraphics[width=0.9\textwidth]{content/source/The groups require non-functional systems.png}
    \caption{Các nhóm yêu cầu phi chức năng của hệ thống}
    \label{fig:non_functional_groups}
\end{figure}

\subsection{Yêu cầu về tính đúng nguồn}

Tính đúng nguồn là yêu cầu quan trọng nhất đối với hệ thống hỏi đáp có sử dụng RAG. Đối với các câu hỏi cần truy xuất tài liệu, câu trả lời của hệ thống phải dựa trên kho tài liệu đã được upload, xử lý và ingest vào hệ thống. Chatbot không nên tự tạo ra câu trả lời vượt quá phạm vi dữ liệu hiện có.

Trong hệ thống, các tài liệu tuyển sinh, điểm chuẩn, chương trình đào tạo, quy chế và biểu mẫu được đưa vào kho RAG. Sau đó, hệ thống xử lý tài liệu thành các đoạn nhỏ, gọi là chunks, để phục vụ quá trình truy xuất. Khi người dùng đặt câu hỏi, RAG service cần tìm các chunks liên quan trước, sau đó mới sinh câu trả lời dựa trên nội dung truy xuất được.

Yêu cầu này giúp hạn chế hiện tượng chatbot trả lời theo suy đoán hoặc tạo ra thông tin không có trong tài liệu. Đặc biệt, với các câu hỏi về điểm chuẩn, học phí, quy chế tuyển sinh hoặc thời gian xét tuyển, hệ thống cần ưu tiên trả lời dựa trên dữ liệu chính thức đã được quản trị viên kiểm soát.

\begin{table}[H]
\centering
\caption{Yêu cầu về tính đúng nguồn}
\label{tab:nfr_source}
\begin{tabular}{|p{2.5cm}|p{11.5cm}|}
\hline
\textbf{Mã yêu cầu} & \textbf{Nội dung} \\ \hline
NFR-01 & Câu trả lời RAG phải dựa trên tài liệu đã được ingest vào hệ thống \\ \hline
NFR-02 & Không sinh câu trả lời vượt quá phạm vi dữ liệu tìm thấy trong kho tài liệu \\ \hline
NFR-03 & Với câu hỏi không tìm thấy nguồn phù hợp, hệ thống cần thông báo không có đủ dữ liệu \\ \hline
NFR-04 & Các tài liệu dùng cho RAG cần được kiểm duyệt trước khi sử dụng \\ \hline
NFR-05 & Câu trả lời về tuyển sinh, học phí, điểm chuẩn và quy chế cần ưu tiên nguồn chính thức \\ \hline
\end{tabular}
\end{table}

Ví dụ, nếu người dùng hỏi: “Điểm chuẩn ngành Công nghệ thông tin năm 2025 là bao nhiêu?”, hệ thống cần truy xuất tài liệu điểm chuẩn hoặc file tuyển sinh đã được xử lý. Nếu không tìm thấy dữ liệu phù hợp, chatbot không nên tự đoán điểm chuẩn.

\subsection{Yêu cầu về khả năng truy vết}

Khả năng truy vết giúp quản trị viên theo dõi được quá trình xử lý tài liệu và kiểm tra lại nguồn dữ liệu khi cần. Mỗi tài liệu trong hệ thống cần có trạng thái xử lý, số lượng chunks, thời điểm cập nhật và mã theo dõi hoặc track id.

Đối với hệ thống RAG, truy vết là yêu cầu quan trọng vì chất lượng trả lời phụ thuộc trực tiếp vào chất lượng tài liệu và trạng thái xử lý của tài liệu. Nếu một tài liệu chưa được xử lý thành công, đang pending hoặc bị failed, hệ thống cần hiển thị rõ để quản trị viên biết và xử lý.

Các thông tin cần truy vết gồm:

\begin{table}[H]
\centering
\caption{Các thông tin cần truy vết}
\label{tab:traceability_info}
\begin{tabular}{|p{5.5cm}|p{8.5cm}|}
\hline
\textbf{Thành phần cần truy vết} & \textbf{Ý nghĩa} \\ \hline
Tên tài liệu & Xác định tài liệu được đưa vào hệ thống \\ \hline
Trạng thái xử lý & Cho biết tài liệu đã processed, pending hay failed \\ \hline
Số chunks & Cho biết tài liệu được chia thành bao nhiêu đoạn \\ \hline
Thời điểm cập nhật & Xác định lần xử lý hoặc cập nhật gần nhất \\ \hline
Track id & Phục vụ theo dõi tiến trình xử lý tài liệu \\ \hline
Dung lượng hoặc độ dài nội dung & Hỗ trợ đánh giá quy mô tài liệu \\ \hline
Nguồn sử dụng & Xác định tài liệu có được dùng cho RAG hay không \\ \hline
\end{tabular}
\end{table}

\begin{table}[H]
\centering
\caption{Yêu cầu về khả năng truy vết}
\label{tab:nfr_traceability}
\begin{tabular}{|p{2.5cm}|p{11.5cm}|}
\hline
\textbf{Mã yêu cầu} & \textbf{Nội dung} \\ \hline
NFR-06 & Mỗi tài liệu cần có trạng thái xử lý rõ ràng \\ \hline
NFR-07 & Hệ thống cần hiển thị số chunks của từng tài liệu \\ \hline
NFR-08 & Hệ thống cần lưu thời điểm cập nhật hoặc xử lý tài liệu \\ \hline
NFR-09 & Hệ thống cần có track id hoặc mã theo dõi để phục vụ kiểm tra pipeline \\ \hline
NFR-10 & Quản trị viên có thể kiểm tra tài liệu nào đang sẵn sàng dùng cho RAG \\ \hline
\end{tabular}
\end{table}

Nhờ khả năng truy vết, khi chatbot trả lời sai hoặc không tìm được dữ liệu, quản trị viên có thể kiểm tra lại tài liệu liên quan đã được upload chưa, đã xử lý thành công chưa và có sinh đủ chunks hay không.

\subsection{Yêu cầu về khả năng mở rộng}

Hệ thống cần có khả năng mở rộng khi số lượng tài liệu, người dùng và câu hỏi tăng lên. Khả năng mở rộng được thể hiện qua việc hệ thống tách các lớp lưu trữ theo vai trò khác nhau, không gom toàn bộ dữ liệu vào một nơi duy nhất.

Theo cấu hình RAG health đã quan sát, hệ thống sử dụng các thành phần lưu trữ như RedisKVStorage, RedisDocStatusStorage, MemgraphStorage và QdrantVectorDBStorage. Cách tổ chức này phù hợp với kiến trúc tách lớp, trong đó mỗi loại lưu trữ đảm nhiệm một vai trò riêng: Redis phục vụ key-value và trạng thái tài liệu, Memgraph phục vụ dữ liệu quan hệ/đồ thị tri thức, còn Qdrant phục vụ lưu trữ vector embedding cho truy xuất ngữ nghĩa.

\begin{table}[H]
\centering
\caption{Các thành phần lưu trữ của hệ thống}
\label{tab:storage_components}
\begin{tabular}{|p{5.5cm}|p{8.5cm}|}
\hline
\textbf{Thành phần} & \textbf{Vai trò} \\ \hline
RedisKVStorage & Lưu dữ liệu key-value, cache hoặc trạng thái xử lý nhanh \\ \hline
RedisDocStatusStorage & Lưu trạng thái xử lý tài liệu \\ \hline
MemgraphStorage & Lưu trữ quan hệ tri thức phục vụ graph context \\ \hline
QdrantVectorDBStorage & Lưu vector embedding phục vụ tìm kiếm semantic \\ \hline
Document Store/Input Directory & Lưu tài liệu gốc hoặc tài liệu đầu vào \\ \hline
\end{tabular}
\end{table}

\begin{table}[H]
\centering
\caption{Yêu cầu về khả năng mở rộng}
\label{tab:nfr_scalability2}
\begin{tabular}{|p{2.5cm}|p{11.5cm}|}
\hline
\textbf{Mã yêu cầu} & \textbf{Nội dung} \\ \hline
NFR-11 & Hệ thống cần hỗ trợ bổ sung thêm tài liệu mới mà không ảnh hưởng dữ liệu cũ \\ \hline
NFR-12 & Kiến trúc lưu trữ cần tách biệt metadata, trạng thái, vector và quan hệ tri thức \\ \hline
NFR-13 & Có thể mở rộng số lượng tài liệu và chunks trong kho RAG \\ \hline
NFR-14 & Có thể thay đổi hoặc nâng cấp thành phần lưu trữ khi quy mô hệ thống tăng \\ \hline
NFR-15 & Có thể mở rộng thêm loại tài liệu như PDF, DOCX, XLSX hoặc dữ liệu văn bản khác \\ \hline
\end{tabular}
\end{table}

Khả năng mở rộng giúp hệ thống không chỉ phục vụ một số tài liệu tuyển sinh ban đầu, mà còn có thể bổ sung thêm nhiều tài liệu về đào tạo, quy chế, biểu mẫu, điểm chuẩn, học phí và chương trình học trong tương lai.

\subsection{Yêu cầu về khả năng vận hành}

Khả năng vận hành thể hiện ở việc hệ thống có thể được theo dõi, kiểm tra và xử lý lỗi trong quá trình chạy thực tế. Đối với hệ thống RAG, quá trình ingest tài liệu có thể gặp lỗi do định dạng file, lỗi trích xuất văn bản, lỗi OCR, lỗi chunking hoặc lỗi sinh embedding. Vì vậy, hệ thống cần có cơ chế theo dõi pipeline và xử lý lại các trường hợp thất bại.

Trong hệ thống hiện tại, phân hệ quản trị tài liệu ngữ cảnh có các thao tác như pipeline status, retry failed, clear cache, refresh và các màn hình thống kê. Những chức năng này giúp quản trị viên dễ dàng theo dõi trạng thái xử lý tài liệu và can thiệp khi cần.

\begin{table}[H]
\centering
\caption{Yêu cầu về khả năng vận hành}
\label{tab:nfr_operability}
\begin{tabular}{|p{2.5cm}|p{11.5cm}|}
\hline
\textbf{Mã yêu cầu} & \textbf{Nội dung} \\ \hline
NFR-16 & Hệ thống cần hiển thị trạng thái pipeline xử lý tài liệu \\ \hline
NFR-17 & Hệ thống cần hỗ trợ xử lý lại các tài liệu bị failed \\ \hline
NFR-18 & Hệ thống cần hỗ trợ làm mới dữ liệu hiển thị trên giao diện quản trị \\ \hline
NFR-19 & Hệ thống cần có chức năng clear cache khi cần cập nhật hoặc xử lý lại dữ liệu \\ \hline
NFR-20 & Hệ thống cần cung cấp các màn hình thống kê để theo dõi hoạt động \\ \hline
NFR-21 & Quản trị viên cần biết tài liệu nào đã xử lý thành công, tài liệu nào đang chờ và tài liệu nào bị lỗi \\ \hline
\end{tabular}
\end{table}

Khả năng vận hành tốt giúp hệ thống phù hợp hơn với môi trường sử dụng thực tế. Khi có tài liệu tuyển sinh mới, quản trị viên có thể upload hoặc scan tài liệu, theo dõi quá trình xử lý, xử lý lại tài liệu lỗi và kiểm tra kết quả sau khi ingest.

\subsection{Yêu cầu về bảo mật và phân quyền}

Bảo mật và phân quyền là yêu cầu bắt buộc đối với hệ thống có phân hệ quản trị. Không phải người dùng nào cũng được phép truy cập vào các chức năng như quản lý tài liệu, upload tài liệu RAG, xem thống kê, quản lý vai trò, quản lý quyền hoặc quản lý tài khoản.

Hệ thống hiện tại có giao diện RBAC, quản lý vai trò, quản lý quyền và quản lý người dùng. Điều này cho phép phân chia quyền truy cập theo vai trò, chẳng hạn người dùng thường chỉ được hỏi đáp, người quản trị nội dung được quản lý tài liệu, còn người quản trị hệ thống có quyền quản lý vai trò, quyền và tài khoản.

Ngoài ra, tài liệu/form có phân biệt public/private trong phần thống kê tài liệu. Điều này cho thấy hệ thống cần kiểm soát phạm vi hiển thị và sử dụng tài liệu, tránh để các tài liệu nội bộ bị truy cập không đúng đối tượng.

\begin{table}[H]
\centering
\caption{Yêu cầu về bảo mật và phân quyền}
\label{tab:nfr_security2}
\begin{tabular}{|p{2.5cm}|p{11.5cm}|}
\hline
\textbf{Mã yêu cầu} & \textbf{Nội dung} \\ \hline
NFR-22 & Người dùng cần được phân quyền theo vai trò \\ \hline
NFR-23 & Chức năng quản trị chỉ dành cho tài khoản có quyền phù hợp \\ \hline
NFR-24 & Hệ thống cần có chức năng quản lý roles và permissions \\ \hline
NFR-25 & Hệ thống cần quản lý danh sách người dùng và trạng thái tài khoản \\ \hline
NFR-26 & Tài liệu/form cần hỗ trợ phân biệt public/private \\ \hline
NFR-27 & Các API quản trị cần kiểm tra quyền trước khi thực hiện thao tác \\ \hline
NFR-28 & Người dùng không có quyền không được phép truy cập chức năng quản lý tài liệu, thống kê hoặc RBAC \\ \hline
\end{tabular}
\end{table}

Yêu cầu này giúp hệ thống đảm bảo an toàn dữ liệu, hạn chế rủi ro chỉnh sửa sai thông tin tuyển sinh và kiểm soát được phạm vi truy cập của từng nhóm người dùng.

\subsection{Bảng tổng hợp yêu cầu phi chức năng}

\begin{table}[H]
\centering
\caption{Tổng hợp yêu cầu phi chức năng}
\label{tab:nfr_summary}
\begin{tabular}{|p{3.2cm}|p{3.3cm}|p{7.5cm}|}
\hline
\textbf{Nhóm yêu cầu} & \textbf{Mã yêu cầu} & \textbf{Nội dung chính} \\ \hline
Tính đúng nguồn & NFR-01 đến NFR-05 & Câu trả lời RAG phải dựa trên tài liệu đã ingest, không trả lời vượt quá dữ liệu \\ \hline
Khả năng truy vết & NFR-06 đến NFR-10 & Mỗi tài liệu có trạng thái xử lý, số chunks, thời điểm cập nhật và track id \\ \hline
Khả năng mở rộng & NFR-11 đến NFR-15 & Tách lớp lưu trữ metadata, trạng thái, đồ thị và vector; hỗ trợ mở rộng tài liệu/chunks \\ \hline
Khả năng vận hành & NFR-16 đến NFR-21 & Có pipeline status, retry failed, clear cache, refresh và các màn hình thống kê \\ \hline
Bảo mật và phân quyền & NFR-22 đến NFR-28 & Có RBAC, roles, permissions, users management và phân biệt public/private \\ \hline
\end{tabular}
\end{table}

\subsection{Nhận xét}

Các yêu cầu phi chức năng trên cho thấy hệ thống KMA Chatbot không chỉ tập trung vào việc trả lời câu hỏi, mà còn chú trọng đến tính tin cậy, khả năng kiểm soát và khả năng vận hành lâu dài. Đặc biệt, với một hệ thống tư vấn tuyển sinh có tích hợp RAG, yêu cầu “đúng nguồn” và “truy vết được tài liệu” là rất quan trọng. Đây là cơ sở để đảm bảo câu trả lời của chatbot có căn cứ, có thể kiểm tra lại và không gây hiểu nhầm cho người dùng.

Bên cạnh đó, việc hệ thống sử dụng các thành phần lưu trữ tách biệt như Redis, Memgraph và Qdrant cho thấy kiến trúc có định hướng mở rộng. Các chức năng quản trị như pipeline status, retry failed, clear cache, thống kê, RBAC và users management giúp hệ thống phù hợp hơn với môi trường triển khai thực tế, nơi dữ liệu thường xuyên thay đổi và cần được kiểm soát bởi nhiều nhóm người dùng khác nhau.
